\subsection{残った有理例外値 $s=-1/3$ の二パラメータ族}

先行する exact elimination の証明書
\texttt{31\_NN\_k1\_IB\_rational\_certificate.json} は、残る変数を
$z_1,z_{11},z_{18},z_{20}$ の四つに帰着する。
残る方程式に
\[
e_1=z_{11}/3-z_{18}z_{20},\qquad
e_2=-z_1/4+z_{11}/12-z_{18}z_{20}/4+3z_{20}^2/4
\]
が含まれ、
\[
z_{11}-3z_{18}z_{20}=3e_1,\qquad
z_1-3z_{20}^2=-4e_2+e_1
\]
なので、パラメータで割らずに
\[
z_{11}=3z_{18}z_{20},\qquad z_1=3z_{20}^2
\]
を得る。これらを代入すると、残りの全方程式は恒等的に零となる。
従って $a=z_{18},b=z_{20}$ を自由変数とする二パラメータ族を得る。

この族の全 Poisson 行列と退化五次式を別添 JSON に保存した。
直接計算により、全 Jacobi 方程式と全必要 $\widetilde E_8$ first-jet 方程式を満たす。
しかし五次式は、全ての $a,b$ に対して
\[
F\in(x_2,x_3)^2.
\]
しかも $x_1^3x_3^2$ の係数は常に $-1/120$ なので $F$ は恒等的に零ではない。
従って全てのパラメータで
\[
V(x_2,x_3)\cong\PP^2\subset\operatorname{Sing}V(F)
\]
となり、正規性に反する。
この例外値には実際の Jacobi 解が存在するが、正規な退化因子を与える点は存在しない。
